% ============================================================
%  Chapter 7 — Robots in Space and Deep Sea
%  Robot World: Meet the Machines That Help Us
% ============================================================

\chapter{Robots in Space and Deep Sea}

\chapteropening{%
  Some places are too hot, too cold, too far, or too deep for humans.
  That is where robots really shine \textemdash{} they go where we cannot.
}
\chapterbadge{Sprocket and Tread: Exploring the Unknown}

% --- Space robots ---
\section*{Robots in Space}

Space is one of the most dangerous environments imaginable: no air, extreme
temperatures, and deadly \term{radiation}. Sending a robot instead of a person
is often safer, cheaper, and sometimes the only option.

\character{Sprocket}{I am inspired by the \term{Mars rovers}! They roll across
  another planet, drill into rocks, and send photos back to Earth
  \textemdash{} all by themselves, because radio signals take up to
  20 minutes to travel between Mars and Earth.}

\subsection*{Mars Rovers}

NASA has sent several rovers to Mars:
\begin{itemize}
  \item \textbf{Sojourner} (1997) \textemdash{} the first, about the size of a microwave.
  \item \textbf{Spirit and Opportunity} (2004) \textemdash{} golf-cart-sized explorers.
  \item \textbf{Curiosity} (2012) \textemdash{} car-sized, still exploring today.
  \item \textbf{Perseverance} (2021) \textemdash{} searching for signs of ancient life.
\end{itemize}

\begin{funfact}
  Perseverance carries a tiny helicopter called Ingenuity \textemdash{}
  the first aircraft ever to fly on another planet!
\end{funfact}


% --- Diagram ---
\begin{diagram}[A Mars rover and its key parts.]
  \begin{tikzpicture}[>=Stealth]
    % Rover body
    \draw[fill=MetalGray!20,draw=SteelBlue,line width=1.2pt,rounded corners=4pt]
      (0,1) rectangle (5,3);
    \node[font=\small\bfseries,align=center] at (2.5,2) {Rover Body\\(computer + instruments)};
    % Wheels
    \foreach \x in {0.7, 2.5, 4.3} {
      \draw[fill=SteelBlue!30,draw=SteelBlue,line width=1pt] (\x,0.5) circle (0.45);
    }
    \node[font=\scriptsize] at (2.5,0) {6 wheels with independent motors};
    % Solar panel
    \draw[fill=SunYellow!30,draw=SunYellow!70!black,line width=1pt]
      (0.5,3) -- (0.2,4) -- (2.3,4) -- (2,3) -- cycle;
    \node[font=\scriptsize\bfseries,above] at (1.25,4) {Solar panel};
    % Camera mast
    \draw[SteelBlue,line width=1.5pt] (3.8,3) -- (3.8,4.2);
    \draw[fill=WarmOrange!20,draw=WarmOrange,line width=1pt,rounded corners=2pt]
      (3.3,4.2) rectangle (4.3,4.7);
    \node[font=\scriptsize\bfseries,above] at (3.8,4.7) {Cameras};
    % Antenna
    \draw[CoralRed,line width=1.2pt] (4.5,3) -- (4.8,3.8);
    \draw[fill=CoralRed!20,draw=CoralRed,line width=0.8pt] (4.8,3.8) circle (0.2);
    \node[font=\scriptsize\bfseries,right] at (5.05,3.8) {Antenna};
    % Robot arm
    \draw[RoboGreen!70!black,line width=1.5pt] (0,2) -- (-0.8,2.5) -- (-1.4,1.8);
    \node[font=\scriptsize\bfseries,left] at (-1.5,1.8) {Drill arm};
  \end{tikzpicture}
\end{diagram}

\sectionrule

\subsection*{Robotic Arms on the Space Station}

The International Space Station has a giant \term{robotic arm} called
\textbf{\term{Canadarm2}}. It catches visiting spacecraft, moves equipment,
and helps astronauts during spacewalks.

\sectionrule

% --- Deep-sea robots ---
\section*{Robots Under the Sea}

The deep ocean is almost as alien as space. At the bottom of the
Mariana Trench \textemdash{} 11 kilometres down \textemdash{} the pressure
would crush a human instantly.

\subsection*{ROVs \textemdash{} Remotely Operated Vehicles}

\term{ROVs} are connected to a ship by a long cable. A human pilot
controls them from the surface using cameras and joysticks.
They repair undersea cables, inspect oil rigs, and film deep-sea creatures.

\subsection*{AUVs \textemdash{} Autonomous Underwater Vehicles}

\term{AUVs} have no cable. They are programmed with a mission, dive down,
collect data, and return on their own. Scientists use them to map the
ocean floor and monitor coral reefs.

\begin{picturethis}
  Imagine swimming to the bottom of 40 stacked swimming pools \textemdash{}
  that is how deep some AUVs can go. It is pitch black, freezing cold,
  and the water presses on every surface with enormous force.
\end{picturethis}

\sectionrule

% --- \term{Rescue robots} ---
\section*{Rescue Robots}

\character{Tread}{When a building collapses or a volcano erupts,
  I go in first. I am built tough \textemdash{} heat, dust, and rubble
  cannot stop me. I search for survivors with cameras and heat sensors.}

Rescue robots:
\begin{itemize}
  \item Crawl through rubble after earthquakes.
  \item Fly drones over wildfire zones to find trapped people.
  \item Carry supplies into flooded areas.
\end{itemize}

\begin{bigidea}
  Robots do not get scared, do not need oxygen, and do not feel pain.
  That makes them perfect for jobs that are too dangerous for humans.
\end{bigidea}

\sectionrule


% --- Comparison table ---
\section*{Comparing Space and Sea Robots}

\begin{center}
\begin{tabularx}{0.9\textwidth}{|X|X|X|}
  \hline
  & \textbf{Space Robots} & \textbf{Deep-Sea Robots} \\
  \hline
  \textbf{Challenge} & No air, radiation, extreme cold/heat & Crushing pressure, darkness \\
  \hline
  \textbf{Power} & Solar panels or nuclear batteries & Battery or cable from ship \\
  \hline
  \textbf{Communication} & Radio (slow --- \term{signal delay}) & Cable (ROV) or acoustic (AUV) \\
  \hline
  \textbf{Movement} & Wheels (rovers) or thrusters & Propellers or fins \\
  \hline
  \textbf{Example} & Perseverance (Mars) & Jason (deep-sea ROV) \\
  \hline
\end{tabularx}
\end{center}

\begin{funfact}
  The deepest-diving robot, \textbf{Nereus}, reached the bottom of the
  Mariana Trench --- 10,902 metres below the surface. That is deeper
  than Mount Everest is tall!
\end{funfact}

\sectionrule

% --- Experiment ---
\begin{experiment}
  \textbf{Simulate a signal delay.}
  \begin{enumerate}
    \item Ask a friend to stand across the playground.
    \item Give them instructions by walkie-talkie, but wait 10 seconds
      before each command (simulating the Mars signal delay).
    \item Try to guide them around cones.
  \end{enumerate}
  \textit{This shows why Mars rovers must be partly autonomous \textemdash{}
  waiting for instructions from Earth is too slow!}
\end{experiment}

\sectionrule


\begin{quickcheck}
  \begin{enumerate}
    \item Why do Mars rovers need to be partly autonomous?
    \item What is the difference between an ROV and an AUV?
    \item Name one dangerous job that rescue robots can do.
  \end{enumerate}
\end{quickcheck}

\sectionrule
% --- New Words ---
\begin{newwords}
  \begin{description}[labelwidth=3.8cm, leftmargin=4.0cm]
    \item[Rover] A wheeled robot designed to travel across a planet's surface.
    \item[ROV] Remotely Operated Vehicle \textemdash{} controlled by a human via cable.
    \item[AUV] Autonomous Underwater Vehicle \textemdash{} navigates on its own.
    \item[Signal delay] The time a radio message takes to travel a long distance.
    \item[Radiation] Invisible energy in space that is harmful to living things.
    \item[Autonomous] Able to operate without human control.
  \end{description}
\end{newwords}
